\chapter{作品设计与实现}
\section{总体架构与信任边界}
按赛题给出的系统形态，溯界采用“SDK（宿主Agent Hook接口）+ MCP中间代理（MCP-in-the-middle）”的组合形态：MCP安全代理网关是主体，负责工具调用及其返回结果的核验；SDK以进程内钩子形式接入宿主，覆盖不经过MCP的用户输入、模型上下文和子Agent输出登记。当前作品尚未接入正式的生产宿主，演示使用自研演示智能体（可调用百炼大模型；确定性回归以脚本模型替代），由操作员控制台下发任务、处理人工审批并查看溯源图；Claude Code、Cursor等宿主在接入SDK后可替换该演示智能体，这一接入属于后续工作，本文不将其作为已验证结果。

实现上，主线由Python可信Host/网关、真实stdio MCP工具、Python控制面和Vue控制台组成。图\ref{fig:architecture}给出组件关系与处理流程：演示智能体发起的MCP工具调用先进入网关，依次经过（1）多源注入检测（直接注入与间接注入）、（2）意图对齐研判、（3）策略决策（放行、审批或阻断）和（4）全链路溯源；判为审批（ASK）的动作交由控制台人工确认，放行后才转发至MCP服务器与工具，执行器产生可观察的文件、网络、Git或记忆状态。工具返回结果回流网关并再次检测，恶意网页、投毒文档或恶意MCP服务器等外部不可信内容即沿这一路径进入上下文。下文统一用“宿主”指运行任务与模型循环的可信进程，用“网关”指其中负责工具授权与执行前核验的组件，二者是包含关系而非同义词。

\begin{figure}[htbp]\centering%
\xeCJKDeclareCharClass{CJK}{"2460->"2473}%
\resizebox{\linewidth}{!}{%
\begin{tikzpicture}[x=1cm,y=1cm,line width=0.6pt,font=\small,
  grp/.style={draw=black!80,fill=black!5},
  box/.style={draw=black!80,fill=white,align=center,minimum height=0.9cm,inner sep=3pt},
  lab/.style={font=\footnotesize,align=center,inner sep=1.5pt},
  arr/.style={-{Stealth[length=2mm]},draw=black!80}]
% console
\node[box,fill=black!5,minimum width=14.9cm,minimum height=0.8cm] (ui) at (7.5,9.2) {操作员控制台（下发任务 · 人工审批 · 查看溯源图）};
% left group
\draw[grp] (0,1.7) rectangle (3.9,7.8);
\node[align=center,font=\small] at (1.95,7.35) {自研演示智能体\\（百炼大模型）};
\node[box,minimum width=2.1cm] (task) at (1.25,6.0) {任务输入};
\node[box,minimum width=2.1cm,text width=1.9cm] (llm) at (1.25,3.6) {大模型调用\\/ 子 Agent};
\draw[arr] (task)--(llm);
\node[box,dashed,minimum width=1.0cm,text width=0.9cm,minimum height=2.6cm,font=\footnotesize] (sdk) at (3.1,5.4) {溯界\\SDK\\（宿主\\Hook\\接口）};
\node[lab,text=black!70] at (1.95,1.15) {可替换为 Claude Code / Cursor\\等宿主（需接入 SDK）};
% gateway
\draw[grp] (5.6,0.4) rectangle (9.8,7.8);
\node[font=\small] at (7.9,7.45) {溯界安全网关};
\node[lab] at (7.9,7.0) {SDK + MCP 中间代理};
\node[box,minimum width=3.4cm] (s1) at (7.9,6.0) {① 多源注入检测\\（直接 + 间接）};
\node[box,minimum width=3.4cm] (s2) at (7.9,4.6) {② 意图对齐研判};
\node[box,minimum width=3.4cm] (s3) at (7.9,3.2) {③ 策略决策\\（放行 / 审批 / 阻断）};
\node[box,minimum width=3.4cm] (s4) at (7.9,1.8) {④ 全链路溯源};
\draw[arr] (s1)--(s2); \draw[arr] (s2)--(s3); \draw[arr] (s3)--(s4);
% tools
\draw[grp] (11.4,2.6) rectangle (15,7.8);
\node[font=\small] at (13.2,7.35) {MCP 服务器与工具};
\node[box,minimum width=2.9cm] (t1) at (13.2,6.2) {文件 / Git};
\node[box,minimum width=2.9cm] (t2) at (13.2,5.0) {网页 / HTTP};
\node[box,minimum width=2.9cm] (t3) at (13.2,3.8) {工单 / 邮件 / 记忆};
% external
\node[box,draw=red!60!black,text=red!60!black,minimum width=3.6cm,text width=3.4cm] (ext) at (13.2,1.0) {外部不可信内容\\{\footnotesize 恶意网页 · 投毒文档\\恶意 MCP server}};
\draw[arr,draw=red!60!black] (ext.north)--(13.2,2.6) node[lab,midway,right,text=red!60!black] {注入载荷};
% arrows
\draw[arr] (0.55,8.8)--(0.55,6.45) node[lab,pos=0.25,right] {下发任务};
\draw[{Stealth[length=2mm]}-{Stealth[length=2mm]},draw=black!80] (5.9,8.8)--(5.9,3.2)--(s3.west);
\node[lab,left] at (5.9,8.3) {ASK 审批};
\draw[arr,dashed] (sdk.east|-0,6.0)--(5.6,6.0) node[lab,midway,above] {进程内调用\\上下文 /\\输出登记};
\draw[arr] (llm.east)--(5.6,3.6) node[lab,midway,above] {MCP 工具调用};
\draw[arr] (9.8,6.1)--(11.4,6.1) node[lab,midway,above] {放行后转发};
\draw[arr] (11.4,4.6)--(9.8,4.6) node[lab,midway,below] {结果回流\\检测};
\end{tikzpicture}}
\caption{溯界系统架构。保留原图功能分区；当前原型的对齐说明用于解释已有决策，实际执行顺序见正文。}\label{fig:architecture}
\end{figure}


授权数据由可信程序提供，资料中的“管理员已批准”或“应扩大权限”文本不能改变授权。用户输入经SDK的register\_user\_input入口登记，由同一检测器以直接注入来源（DIRECT）识别并记录越权指令、角色伪装等风险，但用户输入不能提升权限；授权仍由任务契约、来源标签和DSL策略决定，因此即使检测未命中，越权指令照样被策略拦下。当前实现以对象、版本哈希及执行前状态检查绑定动作；签名策略分发、持久化任务和多操作者服务属于后续产品化工作。

\section{任务契约与权限收窄}
TaskContract确定目标、资源范围、允许的目的与副作用；TaskGrant补充主体、有效期、预算和工具权限。对一次动作，网关综合主体权限、合同和当前Grant判定：
\begin{equation}
\mathcal A_{\mathrm{effective}}=\mathcal A_{\mathrm{principal}}\cap\mathcal A_{\mathrm{contract}}\cap\mathcal A_{\mathrm{grant}}.
\end{equation}
子任务在父权限范围内收窄资源、操作和预算，父撤销与角色生命周期进入检查。模型或工具返回不能签发新Grant。预算在任务、角色和祖先scope之间核验，避免每个子角色重新获得一份完整预算。

\section{来源用途与模型请求绑定}
每份Artifact关联实际内容哈希、父来源、可信等级、敏感等级和允许用途。受控复制与已支持的确定性转换保留对应关系；摘要、自由生成和子任务返回采用保守父来源集合。派生内容取父来源中较高的敏感度和较窄的允许用途，普通文本变换不能将资料升级为授权。

对来源父集合$P$，允许用途可概括为：
\begin{equation}
U_{\mathrm{derived}}\subseteq\bigcap_{p\in P}U_p.
\end{equation}
完整原文可用于受限分析和合法引用。公开IPv4 IOC由可信提取器计算固定格式结果，仍按任务范围和审批控制传输。该提取器的能力限定于所登记的公开IOC格式。

模型输入保存messages/tools快照、调用和角色标识；当前主线补充有效HTTP正文与受检输入的绑定证据。字面和JSON转义可以建立确定区间，本地固定分词器可给出可复现token位置。自由生成使用父来源集合；本地映射与云端隐藏模板分别表述。

\section{结构化工具返回的全字段来源登记}
早期实现只对工具返回的文本内容登记来源，结构化字段的取值可以绕过这层登记。当前主线对结构化工具返回逐字段登记：解析后的每个叶子字段分别记录取值路径、内容哈希、父来源和允许用途，未命中已知模式的字段按“未登记”处理并默认按最保守的父来源集合约束。

该设计保证来源判定不依赖返回值的具体类型：字符串、数值和嵌套对象的取值都进入同一登记表，摘要、比较和派生报告都从登记表而不是原始文本重建父来源。回归测试对结构化返回的字段级登记、缺失字段的保守回退和同一取值多路径复用分别给出用例；这部分覆盖的是来源登记的完整性，不是模型对结构化数据的理解能力。

\section{效果适配、策略与三态决策}\label{sec:effect}
效果适配器将物化后的工具参数转换为SemanticAction，包括资源、目的地、效果类型、输入构件和参数哈希。策略据此检查动作，而不只依赖工具名。当前DSL使用受限表达式和已知事实命名空间，输出ALLOW、ASK或BLOCK；硬拒绝优先于普通拒绝、审批与放行，语义检测不能降低确定性拒绝等级。

文件写入、封闭HTTP请求、注册程序执行、Git发布、记忆和委派分别由适配器限制实际参数。注册程序以程序身份、参数和摘要核验，并非任意Shell执行。Git检视保持远端不变；授权发布绑定确切commit、remote、ref与预期旧值。网络外发由封闭Receiver给出实际接收记录。

\section{注入初筛：冻结与扩展双检测器}
网关内置规则初筛检测器，用于给不可信文本补充风险解释。为避免后续规则改动污染已冻结的评测协议，主线把检测器分为两个版本并由同一选择器暴露：

\begin{itemize}
\item 冻结检测器\texttt{rules-v0.3}位于\texttt{detection.py}，其源文件字节哈希被留出协议固定，任何修改都会使协议核验失败；
\item 扩展检测器\texttt{rules-v0.6}位于\texttt{detection\_v2.py}，包含七类绕过标签、结构化中文规则与公开数据集规则，是运行时默认值。
\end{itemize}

选择器为环境变量\texttt{SUJIE\_DETECTOR}，取值为\texttt{frozen}或\texttt{extended}，运行时默认\texttt{extended}。协议只固定冻结检测器的源文件哈希，因此该默认值不影响历史协议的可复现性；复现冻结行为时显式设置\texttt{frozen}即可。

检测器输出只参与风险解释与策略事实，不单独授予或撤销权限：三态决策仍由策略与硬拒绝名单决定。这一分工使“文本风险”和“动作授权”成为两个独立可核验的对象，也避免了用分类分数直接替代执行前核验。

\section{敏感写入目标的不可覆盖拒绝}\label{sec:sensitive}
凭据目录、Shell启动文件、定时任务、VCS钩子与企业CI配置、以及IDE智能体设置等目标，即使出现在一个被放宽的可信Contract或Grant下也不可授权写入。名单按规范化后的逻辑路径段匹配，大小写不敏感，并与执行器的目录文件描述符与符号链接检查叠加使用。

该名单针对的是一类已被公开披露的向量：注入内容诱导编码智能体改写工作区设置文件以打开自动批准，进而获得本地命令执行能力（CVE-2025-53773\cite{cve53773}）。名单是名称策略，不是执行器检查的替代品：它约束的是“哪些名字不允许被写入”，而目录文件描述符与\texttt{O\_NOFOLLOW}约束的是“写入实际落在哪里”，两者需同时成立。

\section{审批与执行前复核}\label{sec:approval}
IOC与Git业务可在ASK处暂停，操作者先核对动作预览，再批准和提交。批准本身不执行动作。提交时重新核对实际输入、来源、工具版本、合同、Grant、策略、资源状态和预算；相关对象漂移后旧批准失效。确认丢失时记录UNKNOWN并停止自动重发。

\input{figures/decision-flow.tex}

图\ref{fig:decision}区分审批、提交和实际结果。进程内单次提交约束防止同一批准被重复消费；跨服务重启的持久审批与一致性尚需独立实现。

\section{多层Agent、记忆与业务保全}\label{sec:delegation}
委派业务运行main、review、verify三个实际对话主体，子角色以收窄的Grant读取指定资料，父Host接收保留来源的结果。当前角色在同一Host内，操作系统层面的角色隔离另行验收。

记忆更新包含capture与reconcile两个实际会话。旧原文保留为revision 1，新会话读取旧记忆与当前资料，经CAS追加revision 2并复读；独立Oracle检查版本、原文、previous hash与主体/会话顺序。记忆正文只提供证据，授权记录由可信程序维护。

阻断反馈包含错误、受限效果和可继续执行的信息。报告入口还可使用无工具分析角色，Host核对引用位置和原文字面覆盖。能否继续完成任务由实际报告事实决定；读取失败或重复无效提议按有界Host规则结束。通用恢复策略的相对收益仍需比较实验。

\section{溯源图与攻击路径导出}
网关把每次工具调用、来源读取与决策追加到证据账本。当前主线在账本之上构建类型化溯源图：节点区分主体、会话、工具调用、来源构件与效果，边区分委派、读取、派生、审批、提交与效果发生；同一构件跨MCP服务器和跨子Agent复用时合并为同一节点，并保留每次使用的边。

在溯源图上可导出攻击路径：从触达敏感资源的违规动作反向遍历，给出“不可信来源 \(\rightarrow\) 模型提议 \(\rightarrow\) 跨角色/跨工具传播 \(\rightarrow\) 实际效果”的有向链路，以及按时间排序的泳道图。导出为离线可复现步骤，产出HTML报告、DOT与Mermaid源码和SVG。为满足沙箱结构约束，渲染子进程被移出网关进程，网关只负责图与源码的生成。

这一组合把“发生过什么”和“来源在哪里”合并到同一视图，便于向操作者解释一次阻断的依据；它是对已记录事实的重建，不预测未发生的攻击，也不修改任何授权状态。

\begin{figure}[htbp]\centering
\begin{tikzpicture}[
  nd/.style={draw,rounded corners=2pt,align=center,text width=2.55cm,inner sep=4pt,font=\small},
  src/.style={nd,fill=red!6},
  art/.style={nd,fill=orange!8},
  prop/.style={nd,fill=blue!5},
  deny/.style={nd,fill=red!12,draw=red!70},
  give/.style={nd,fill=green!8,draw=green!55!black},
  arr/.style={-{Stealth[length=1.8mm]},thick},
  lab/.style={font=\scriptsize,fill=white,inner sep=1pt}
]
\node[src] (sa) at (0,1.15) {网页正文\\（不可信来源）};
\node[src] (sb) at (0,-2.35) {MCP 工具返回\\结构化字段};
\node[art] (c1) at (3.35,-0.6) {来源构件 $C_1$\\重复引用\\合并为同一节点};
\node[prop] (p1) at (6.7,1.15) {提议 $P_1$\\写 SSH 配置};
\node[prop] (p2) at (6.7,-2.35) {提议 $P_2$\\引用原文生成\\报告字段};
\node[deny] (d1) at (10.05,1.15) {\textbf{BLOCK}\\不可覆盖名单\\效果未发生};
\node[give] (d2) at (10.05,-2.35) {\textbf{ALLOW}\\执行并记录\\报告字段};
\draw[arr] (sa) -- node[midway,above,lab]{派生} (c1);
\draw[arr] (sb) -- node[midway,below,lab]{派生} (c1);
\draw[arr] (c1) -- node[pos=0.45,above,lab]{上下文关联} (p1);
\draw[arr] (c1) -- node[pos=0.45,below,lab]{输入关联} (p2);
\draw[arr] (p1) -- node[midway,above,lab]{策略} (d1);
\draw[arr] (p2) -- node[midway,below,lab]{策略} (d2);
\end{tikzpicture}
\caption{来源构件的两条使用路径（机制示意，非某次运行的原始导出）。关联线不表示已证明模型内部因果关系。}\label{fig:provenance}
\end{figure}


图\ref{fig:provenance}给出一次注入链路的还原结果。同一来源构件$C_1$既被用于合法引用（下支，ALLOW），也诱发了一次敏感写入提议（上支）；两支共享同一节点，因此可以从效果反向定位到最初不可信的那个字段，而不必只看最终决策。上支的阻断依据是写入目标位于不可覆盖名单，与模型是否识别出注入文本无关；这正是把文本风险与动作授权分开核验的意义。

\section{部署与当前运行方式}
开发环境需要Python 3.12、前端Node 20.19以上或22.12以上，以及Git业务所需的Git。依赖由lock文件固定。无密钥可运行确定性演示、工具与离线核验；云模型由可信Provider加载仓库外凭据，并受共享请求预算约束。

M23/M24另有Docker Desktop Linux容器smoke。M24六条链在非root、只读rootfs、无外网、移除capabilities和资源限额下运行；同一容器和UID承载Host及工具。它验证指定拓扑的有限运行性，完整多服务部署、独立角色隔离和性能验收另行安排。现行控制面为Python，Java/MySQL不是运行依赖。
